\section{پاسخ سوال ۹}
\begin{stepbox}
\field{نمایش ماتریس با روش Run Length Coding (RLC)}{
برای فشرده‌سازی بلوک‌های ۸×۸ در تصویر (مشابه استاندارد JPEG)، به دلیل وجود مقادیر صفر فراوان در فرکانس‌های بالا (عموماً در گوشه پایین-راست جدول)، از روش \textbf{کدگذاری طول-اجرا (RLC)} استفاده می‌کنیم. اما پیش از آن، باید مقادیر را با استفاده از \textbf{پیمایش زیگ‌زاگ (Zig-Zag Scan)} استخراج کنیم تا صفرهای متوالی در کنار هم قرار بگیرند و طول رشته‌های صفر بهینه‌تر شود.

\textbf{گام اول: استخراج آرایه با پیمایش زیگ‌زاگ}
با حرکت مورب از فرکانس‌های پایین (گوشه بالا-چپ) به سمت فرکانس‌های بالا (پایین-راست)، آرایه یک‌بعدی زیر به دست می‌آید:
\begin{align*}
\text{ZigZag} = [ & 11, 0, -3, 0, 0, 1, 1, 0, -1, 2, -2, 1, 0, -2, 1, 1, 1, 1, -1, \\
& 0, 0, 1, -2, 1, 0, -1, 0, 0, 0, 0, 0, 1, 0, -2, 2, -1, -1, 1, \\
& 0, 0, -1, 0, 0, 0, 0, 0, 0, 0, 0, 0, -1, 0, 0, \dots, 0 ]
\end{align*}

\textbf{گام دوم: اعمال کدگذاری Run-Length}
در این الگوریتم، به جای ذخیره‌ی خام و پیوسته اعداد، هر درایه غیرصفر را به همراه تعداد صفرهای متوالیِ پیش از آن، به شکل یک زوج‌مرتب $(R, V)$ کدگذاری می‌کنیم:
\begin{itemize}
    \item \textbf{$R$ (\lr{Run}):} تعداد صفرهای متوالی قبل از رسیدن به عدد.
    \item \textbf{$V$ (\lr{Value}):} مقدار خود آن درایه غیرصفر.
\end{itemize}
در نهایت، زمانی که تا انتهای بلوک فقط مقدار صفر باقی مانده باشد، از نشانگر ویژه \textbf{\lr{EOB (End of Block)}} برای اعلام پایان بلوک استفاده می‌کنیم.

\vspace{0.4cm}
نتیجه نهایی کدگذاری RLC برای این جدول به صورت جدول‌بندی‌شده و مرتب به شکل زیر است:
\vspace{0.2cm}
\begin{center}
\renewcommand{\arraystretch}{1.5}
\begin{tabular}{c c c c c}
    $(0, 11)$ & $(1, -3)$ & $(2, 1)$  & $(0, 1)$  & $(1, -1)$ \\
    $(0, 2)$  & $(0, -2)$ & $(0, 1)$  & $(1, -2)$ & $(0, 1)$  \\
    $(0, 1)$  & $(0, 1)$  & $(0, 1)$  & $(0, -1)$ & $(2, 1)$  \\
    $(0, -2)$ & $(0, 1)$  & $(1, -1)$ & $(5, 1)$  & $(1, -2)$ \\
    $(0, 2)$  & $(0, -1)$ & $(0, -1)$ & $(0, 1)$  & $(2, -1)$ \\
    $(9, -1)$ & \textbf{EOB} & & & \\
\end{tabular}
\end{center}
\vspace{0.2cm}

\begin{tipbox}[اهمیت EOB در فشرده‌سازی]
همان‌طور که مشاهده می‌کنید، ۱۳ درایه انتهایی بلوک تماماً صفر هستند. با به کارگیری نماد \lr{EOB}، از ذخیره‌سازی این صفرها جلوگیری شده و حجم داده‌ها به شکل چشم‌گیری کاهش می‌یابد. این ترکیب (تبدیل کسینوسی + زیگ‌زاگ + RLC) هسته اصلی الگوریتم فشرده‌سازی تصویر \lr{JPEG} را تشکیل می‌دهد.
\end{tipbox}
}
\end{stepbox}
